% ============================================================
%  COVER LETTER TEMPLATE
% ------------------------------------------------------------
%  To use: copy this file into an application folder as
%  cover-letter.tex, then fill in the bracketed fields.
%    cp shared/cover-letter-template.tex applications/<company>/cover-letter.tex
%  build.sh auto-detects it and produces
%    dist/<YourName>-CoverLetter-<Company>.pdf
%
%  Name, contact line and region phone come from config.json via
%  shared/identity.tex, same as the resume — nothing to edit here.
%
%  A NOTE ON WHO WRITES THIS. /resume-build will not draft a cover
%  letter for you, by design. A cover letter is the one document
%  that is entirely voice and motive, and a generated one reads
%  like a generated one. The skill flags that a letter is required
%  and leaves it to you.
% ============================================================
\def\TargetRegion{uk}          % override per application if needed
\documentclass[a4paper,11pt]{article}
\usepackage[a4paper,margin=2cm]{geometry}
\usepackage{parskip}
\usepackage{hyperref}
\pagenumbering{gobble}

\input{identity.tex}

\begin{document}

\HeaderBlock

\vspace{1em}
\noindent [Date] \\
\noindent [Hiring Manager / Recruiting Team] \\
\noindent [Company], [Location]

\vspace{1em}
\noindent Dear [Hiring Manager],

\vspace{0.5em}
[Opening: the role you are applying for, and one sentence on why this firm
specifically. If that sentence would work for any firm in the sector, it is not
finished yet.]

[Body: one or two paragraphs connecting a specific thing you have done to a
specific thing the role needs. Name the JD line you are answering. Do not
restate the résumé — the reader has it.]

[Closing: availability and thanks. Short.]

\vspace{0.5em}
\noindent Sincerely, \\
\noindent \CandidateName

\end{document}
